\documentclass[10pt]{article}
\usepackage{template}
\usepackage{amsmath,amssymb,microtype,tikz,graphicx}
\usetikzlibrary{arrows.meta,positioning}
\setlength{\emergencystretch}{2em}
\SetKeys{date=2026-10-20,vol=007,title={每日有趣新闻｜经济、火星与数学},author={ChatGPT 编辑部},note={\small 2026 年 10 月 20 日·第 007 期｜新闻检索截至北京时间 10 月 20 日上午。重点核查 10 月 19 日公开资料与今晨航天发射。原始统计、媒体分析与预印本主张分别标注。}}
\renewcommand{\newsitem}[3]{\par\medskip{\large\bfseries\color{newsblue}#1\par}{\footnotesize\color{newsgray}#2\par}\smallskip\color{black}#3\par\medskip{\color{newsgray}\rule{\linewidth}{0.25pt}}\par}
\newcommand{\newslead}[3]{\par\medskip{\LARGE\bfseries\color{newsblue}#1\par}\vspace{1mm}{\small\color{newsgray}#2\par}\smallskip\color{black}#3\par\medskip{\color{newsgray}\rule{\linewidth}{0.45pt}}\par}
\begin{document}
\begin{documents}
\newslead{中国前三季度经济增长4.6\%：出口与高技术制造拉动增长，投资和消费为何不同步？}{宏观经济头条｜国家统计局 10 月 19 日数据；路透社与美联社交叉核对}{
十月十九日，国家统计局公布二〇二六年前三季度国民经济数据。国内生产总值为一百零七万一千零七十一亿元，按不变价格计算同比增长百分之四点六；三季度单季同比增长百分之四点五，较二季度的百分之四点三回升，但低于一季度的百分之五。三季度环比增长百分之一点二。这组数字说明经济总量仍在扩张，但不能单凭一个总量增速判断企业利润、居民消费和房地产投资是否同时改善。宏观统计的关键恰恰在于不同部门出现了方向不一致的变化。\cite{NBS2026,ReutersChina,APChinaGDP}

\textbf{生产端的扩张相当明显。}前三季度规模以上工业增加值同比增长百分之五点三，其中装备制造业增长百分之十点二，高技术制造业增长百分之十四点三。国家统计局列举的三维打印设备、锂离子电池和工业机器人产量分别增长百分之六十二点五、三十九点五和二十九点一。工业增加值衡量的是生产活动创造的价值，而具体产品产量衡量件数或物理单位；设备产量翻倍并不意味着整个工业产值同步翻倍。高技术制造业占比、行业价格与库存变化，也会影响这些数字最终对应的利润和就业。

对外贸易提供了另一股动力。前三季度以人民币计价的货物进出口总额同比增长百分之十八点八，其中出口增长百分之十六点四，进口增长百分之二十二点三；机电产品出口增长百分之二十三点七。路透社的采访指出，一些企业因为国内竞争激烈而转向海外市场，外需增长能够弥补国内订单不足，却也提高了企业对国际市场和贸易政策的敏感性。报道中的企业案例是个别受访者经历，不能直接推断全部企业均如此；但出口强、内需弱的宏观组合，确实值得继续检验。\cite{NBS2026,ReutersChina}

\begin{center}
\begin{tikzpicture}[x=0.22cm,y=0.33cm]
\draw[->,gray] (-8,0)--(20,0);
\foreach \x in {-5,0,5,10} {\draw[gray!45] (\x,-.18)--(\x,9.3); \node[font=\tiny,gray,below] at (\x,-.18) {\x\%};}
\foreach \label/\val/\y/\col in {工业增加值/5.3/8/newsblue,高技术制造/14.3/7/newsblue,社会消费品零售/1.1/6/newsblue,固定资产投资/-7.2/5/red!70!black,货物进出口/18.8/4/newsblue} {
\pgfmathsetmacro{\start}{min(0,\val)}
\pgfmathsetmacro{\width}{abs(\val)}
\fill[\col!65] (\start,\y-.24) rectangle ++(\width,.49);
\node[font=\tiny,anchor=east] at (-8.5,\y) {\label};
\node[font=\tiny,anchor=west] at (\val+.3,\y) {\val\%};
}
\end{tikzpicture}
\par{\scriptsize 图1　2026年前三季度部分指标同比增速（\%）。工业、高技术制造业为规模以上工业增加值；房地产开发投资同比下降20.6\%，未列入坐标；投资与贸易为各自公布口径。各指标不可直接相加。数据：国家统计局，2026-10-19；本刊绘制。}
\end{center}

\textbf{消费并非没有增长，而是增长幅度与生产端存在差距。}前三季度社会消费品零售总额同比增长百分之一点一，九月份同样为百分之一点一。与此同时，服务零售额增长百分之五，网上商品和服务零售额增长百分之四点四。零售总额主要统计商品和餐饮消费，服务零售采用另一套口径，不能把百分之五与百分之一点一直接相加得到所谓“总消费增速”。国家统计局同时公布，全国居民人均可支配收入实际增长百分之四点四，人均消费支出实际增长百分之二点七。收入与支出增速差异可能涉及储蓄倾向、家庭负债、消费结构等多种机制，不能只凭两个增长率认定某个单一原因。\cite{NBS2026}

投资数据则更加鲜明。前三季度固定资产投资（不含农户）同比下降百分之七点二，房地产开发投资下降百分之二十点六，民间投资下降百分之十点四。房地产新建商品房销售面积下降百分之十三点二，销售额下降百分之十三点九。固定资产投资是一定时期内完成的投资工作量，不等同于上市公司资本开支，也不等同于所有资本形成；在统计口径经过修订的情况下，更应避免机械拼接历年数据。高技术产业投资同期增长百分之四点九，知识产权产品投资增长百分之九点三，说明投资总量下行与部分新兴领域扩张可以同时发生。\cite{NBS2026,NBSInvest}

\textbf{为什么增长率不能简单相减？}有人可能把工业增长百分之五点三、零售增长百分之一点一，直接相减得到“生产过剩百分之四点二”。这是错误的：两者覆盖对象、价格处理、统计基期和单位均不同。工业增加值按可比价格计算，零售总额以商品销售额为主要基础；还有出口、政府消费、库存变化、服务业与进口等环节。它们可以共同提示结构失衡，却不能构成一条简单的会计恒等式。

官方通报强调新动能和就业物价总体稳定，同时承认国内需求不足仍较突出；路透社采访的经济学家则更强调出口依赖和房地产低迷的中期风险。两种叙述关注的是同一数据的不同侧面。后续观察应包括居民实际消费、企业利润、民间投资、地方财政、贸易摩擦和就业质量，而不是只盯住全年能否完成百分之四点五至五的增长目标。\cite{NBS2026,ReutersChina,APChinaGDP}
}

\newslead{凌晨四点四十一分，MMX 启程：日本为何要从火星的卫星带回十克沙土？}{航天深度｜JAXA 10 月 20 日发射结果；NASA 任务说明}{
日本宇宙航空研究开发机构（JAXA）十月二十日发布正式结果：当地时间凌晨四时四十一分零三秒，H3 运载火箭第十号机从种子岛宇宙中心升空；约一小时十一分十三秒后，搭载的火星卫星探测器 MMX 正常分离并进入预定轨道。这里的“成功”特指本次火箭发射与航天器分离，绝不意味着采样、火星抵达或地球返回已经完成。日本当地媒体和美联社也报道了这次发射。\cite{JAXALaunch,APMMX}

MMX 是 Martian Moons eXploration 的缩写，即火星卫星探测计划。火星有两颗形状不规则的小卫星：较靠近火星的火卫一（Phobos）和较远的火卫二（Deimos）。任务计划考察两者的地形、重力和组成，并尝试从火卫一表面采集至少约十克样品，最终在二〇三一年左右送回地球。十克看似很少，但如果样品来自已知位置、经过严格污染控制，并能在实验室用高精度质谱、显微分析与同位素方法反复测量，其科学价值远大于一张模糊的远程照片。\cite{JAXAMMX,NASAMMX}

\begin{center}
\begin{tikzpicture}[node distance=1.1mm, every node/.style={font=\scriptsize}, box/.style={draw=newsblue,rounded corners=1pt,fill=newsblue!5,text width=3.25cm,align=center,minimum height=.66cm}]
\node[box] (a) {2026-10-20　地球发射\quad 已确认};
\node[box,below=of a] (b) {飞往火星、探测火卫一／火卫二\quad 任务计划};
\node[box,below=of b] (c) {火卫一取样、封存、返回\quad 尚未完成};
\node[box,below=of c] (d) {2031 年左右地球接收样品\quad 计划目标};
\foreach \p/\q in {a/b,b/c,c/d} {\draw[-{Stealth[length=1.7mm]},newsblue] (\p)--(\q);}
\end{tikzpicture}
\par{\scriptsize 图3　MMX 任务时间线。发射与分离据 JAXA 10 月 20 日公告确认；后续节点为任务目标，不代表已经完成。本刊绘制。}
\end{center}

\textbf{科学问题的核心是起源。}一种假说认为，火卫一和火卫二原本是小行星，后来被火星引力俘获；另一种假说认为，早期火星受到大型天体撞击，抛出的碎屑逐渐聚集形成卫星。两种机制对岩石化学组成、同位素比例、颗粒结构和轨道演化的预期并不完全相同。火卫一外观黝黑、表面坑洼，但远程观测不足以唯一确定其成因；因此，样品返回是检验假说的一种关键方法。即使最终结果不能在两个假说之间做出简单二选一，也可能迫使研究者提出更复杂的混合形成模型。

这不是第一次从小天体取样。日本此前的隼鸟系列已完成小行星样品返回，但火星卫星的动力学环境、通信时延和航行距离带来了不同挑战。火卫一引力很弱，航天器需要谨慎接近、选择安全着陆点、接触取样，再确保返回器能够长期保存材料。任务还计划投放由法国和德国参与研制的小型巡视器，协助研究地表条件。NASA 提供的 MEGANE 仪器将利用伽马射线和中子测量元素组成；它所提供的是遥感元素信息，不等于对样品的实验室化学分析。\cite{NASAMMX,JAXAMMX}

\textbf{任务时间表是一连串条件，而非已经兑现的结果。}发射成功只是第一关。航天器还需要完成深空航行、火星附近轨道操作、近距离测绘、样品获取与返航，任何环节都可能改变时间表。对比“计划返回十克”与“已成功获得十克”尤其重要；目前能够确认的只有发射、分离和预定轨道注入。后续报道应优先追踪 JAXA 的飞行状态、采样位置选择、载荷观测结果及返回舱实际状态，而不应在新闻标题中提前宣布火星之谜已经破解。\cite{JAXALaunch,JAXAMMX}
}

\newsitem{数学拾趣｜在棋盘上抢四个角：一个胜负分界竟然是 $n=5$}{组合数学与 Lean｜Sean Reid 的 arXiv 预印本；10 月 19 日进入新稿列表}{
两人轮流占据一张 $n\times n$ 方格棋盘中的空格，谁先占有某个边与棋盘边平行的正方形的四个顶点，谁就获胜；棋盘填满而无人达成则平局。这种规则比五子棋简单，但“先手是否一定能赢”不是靠直觉就能判断的。研究者 Sean Reid 在预印本《The square achievement game on $n\times n$ boards》中提出并证明：先手存在必胜策略，当且仅当 $n\ge5$。论文初稿提交于十月十四日，出现在十月十九日 arXiv 数学新稿列表，不能把它误记为十九日才完成的研究。\cite{ReidSquare,ArxivNew}

\begin{center}
\begin{tikzpicture}[scale=.39]
\foreach \x in {0,...,4} {\foreach \y in {0,...,4} {\draw[gray!50] (\x,\y) rectangle ++(1,1);}}
\foreach \x/\y in {0/0,0/2,2/0,2/2} {\fill[newsblue!80] (\x+.5,\y+.5) circle (.31);}
\draw[very thick,newsblue] (.5,.5) rectangle (2.5,2.5);
\node[font=\small,anchor=west] at (6,4.3) {同一玩家占有四角};
\node[font=\small,anchor=west] at (6,3.0) {四角构成轴平行正方形};
\node[font=\small,anchor=west] at (6,1.7) {即满足本游戏获胜条件};
\node[font=\scriptsize,gray,anchor=west] at (6,.4) {本刊示意，非论文棋谱};
\end{tikzpicture}
\par{\scriptsize 图2　正方形达成游戏规则示意。四枚蓝子属于同一玩家，四角构成一个边与棋盘边平行的正方形；本刊依据 Reid（2026）定义绘制。}
\end{center}

这个结论的量词值得细读。所谓“必胜策略”，并不是指先手随便下就能赢，而是存在一套决策规则，使得无论后手如何合法回应，最终先手都能先占齐四角。把每个局面视为一个节点，每个合法落子视为一条边，就得到一棵博弈树。若要证明某个策略必胜，需要说明对方的\emph{每一种}回应，都有继续获胜的路线；只展示一盘漂亮的胜局是不够的。

为什么分界恰好出现在五？小棋盘可用有限状态搜索穷举所有合法回应，但随着 $n$ 增大，可能局面数呈指数级增长。作者对 $5\times5$ 棋盘采用专门的有限表验证；对于 $n\ge6$，则围绕一种“十字形”布局构造统一的强迫性走法，避免为每个更大的棋盘分别列出庞大策略表。论文报告对 $6$ 至 $17$ 的棋盘完成检查，并借助 $23\times23$ 的一次验证与模拟定理推广到更大的棋盘。这里“模拟”不是用有限次随机对局猜测无限情况，而是建立从大棋盘策略行为到已验证结构的数学映射。\cite{ReidSquare}

尤其值得注意的是，作者提供了 Lean 4 形式化文件和相关程序，声称关键策略定理已被形式化。Lean 内核能够检查编码后的定理是否由允许的规则推出，但它不能自动保证人类描述的游戏规则与编码完全一致，也不能代替独立研究者检查所有计算依赖。本刊核对了论文摘要与公开的代码目录，但未在本地重新编译全部形式化工程；因此把结论表述为“作者证明并提供形式化”，而非“本刊独立验证”。这类既有清楚规则又有可运行验证材料的研究，尤其适合数学与编程爱好者进一步复现。\cite{ReidSquare}
}

\newsitem{美国财政赤字约两万亿美元：债务总额、市场持有债务与利息负担有何区别？}{国际经济｜路透社 10 月 19 日分析，援引美国财政部财年数据}{
美国联邦政府在九月三十日结束的二〇二六财年录得约两万亿美元赤字，路透社援引财政部报告称，比上一财年扩大约百分之十二。美国政府债务总额在八月跨过四十万亿美元，其中一部分是政府账户之间的债权债务；公众及市场持有的国债规模约三十二万亿美元。讨论债务占 GDP 比例时，必须说明采用的是总债务还是公众持有债务，否则同一时期可能出现看似互相矛盾的比例。\cite{ReutersUSDebt}

赤字是某一财年支出超过收入的\emph{流量}，债务是历年借款累积形成的\emph{存量}。一个国家即使当年财政赤字缩小，债务总额也可能继续增加；而债务利息又会反过来成为下一年的支出。路透社称，二〇二六财年债务服务成本达到约一点三六万亿美元，赤字约占 GDP 的百分之六点二。债券收益率上升会逐渐提高新债和到期再融资债务的成本，但并不意味着全部存量国债在一夜之间同时按新利率计息。

美国政府强调经济增长与关税收入可能改善财政状况，部分经济学家则担忧高利息和长期刚性支出使赤字难以自然收敛。路透社同时指出，市场并未出现可直接证明美国信用即将崩溃的异常期限溢价。因而，报道既不能把高债务简单等同于马上违约，也不能因为美元仍是国际主要储备货币就断言财政约束不存在。债务结构、期限分布、真实利率、经济增长和未来税收政策，才是理解财政可持续性的关键变量。\cite{ReutersUSDebt}
}

\newsitem{法国骚乱的金融后果：债券收益率为什么比股市下跌更值得观察？}{国际政治与市场｜10 月 19 日路透社；接续第 006 期法国抗议报道}{
法国十月十七日的大规模抗议与部分城市的暴力事件，在十九日交易中进一步传导至金融市场。路透社报道，法国十年期国债收益率当日上升约十三个基点至百分之四点九附近，法国 CAC 40 股指下跌约百分之一点五，部分大型银行股跌幅约百分之三。市场价格随时变动，以上均为报道时点的行情，不是全天最终结算值。\cite{ReutersFranceMarket}

国债收益率与债券价格大体反向变化。投资者担心未来财政预算难以通过或债务风险上升时，可能要求更高收益率才愿意持有债券。十九日法国与德国十年期国债收益率之差约一百四十六个基点，即约一点四六个百分点；它反映市场对两国借款条件的相对定价，却不能被直接解释为法国的违约概率。银行持有债券、依赖市场融资，因此主权利差扩大也可能影响金融机构估值。

抗议者的民生诉求、街头暴力、政府对集会的应对和市场对财政预算的担忧，是相互关联但不能混为一谈的四个问题。路透社报道政府要求部分地方考虑宵禁措施，法国总理并未排除紧急状态选项；这些属于政府当时的政策讨论，不等于全国紧急状态已经生效。后续关键在于预算谈判、独立执法调查和具体经济政策，而不是用股市某天的跌幅给整场社会运动下结论。\cite{ReutersFranceMarket}
}

\newsitem{中美俄三方军控会谈：外交意向尚不等于会议已经确定}{外交观察｜路透社 10 月 19 日报道}{
十月十九日，美国总统特朗普表示愿意考虑下个月与中国国家主席习近平、俄罗斯总统普京进行三方会谈，议题可能包括核军控。路透社报道，这一设想可能与中国举行的亚太经合组织会议时间相近，但截至报道发布时，三方尚未共同宣布正式会议安排、议程或核裁军谈判框架。总统表达意愿、外交部门开始接触和三方签署协议，分别是完全不同的外交阶段。\cite{ReutersTrilateral}

三国核武库规模、部署结构与既有军控义务并不相同。任何多边安排都需要讨论可核查性、现有核力量、导弹与其他投送工具、地区安全环境等问题；如果没有这些细节，仅以“愿意谈”判断核风险已经降低，证据不足。中俄美三方会谈还可能与俄乌战争、贸易和地区安全议题交织，但具体关联必须等待官方议程或可核实谈判结果。\cite{ReutersTrilateral}
}

\newsitem{数学速览｜一篇用 Lean 形式化有损压缩理论的预印本}{信息论与形式化｜arXiv 10 月 19 日数学新稿列表}{
Fubin Yan 等作者在预印本《Aligning Information Theory and Formal Verification》中介绍了对率失真理论核心组成部分的 Lean 形式化工作。率失真理论研究的是允许一定信息损失时，压缩率能够低到什么程度；其核心对象包括失真度量、可达性和渐近编码界。作者强调，教科书的数学定义转写为证明助手可检查的对象时，必须仔细处理概率、极限和量词。论文声称提供可复用的定义与引理，并公开验证脚本；本刊尚未独立运行其工程，因此不将作者的形式化声明视作已经完成的第三方复核。\cite{YanRate}
}

\newsitem{俄罗斯财政部长兼任副总理：战时预算压力进入内阁权力结构}{俄乌战争与经济政策｜路透社 10 月 19 日报道}{
俄罗斯总统普京十月十九日调整政府人事：财政部长安东·西卢安诺夫将兼任副总理，原财政部副部长伊琳娜·奥克拉德尼科娃接任政府办公厅主任。克里姆林宫表示，调整意在强化宏观经济稳定和财政纪律。西卢安诺夫自二〇一一年起任财政部长，长期负责预算和税收事务；其职务扩张显示财政协调在当前政策中的地位上升，但不能仅凭任命推断俄罗斯将立即改变对乌军事政策。\cite{ReutersRussiaCabinet}

路透社援引正在审议的预算修订指出，俄政府把二〇二六年赤字预期从 GDP 的百分之一点六提高到百分之三点二，年度支出预计增加百分之十三点二，至四十八点六万亿卢布；二〇二七年国防支出相较原计划拟上调约百分之二十七。这里的二〇二七年增幅是\emph{相对原预算计划}，并非国防开支同比增长百分之二十七，不能偷换比较基准。财政赤字扩大、利率偏高和税收压力增加，会使预算决策面临取舍：维持战争支出、保障社会支出和限制物价压力，可能要求不同部门接受不同程度的约束。

财政部长同时进入副总理层级，可能增强对其他部门预算执行的协调权；但实际效果还取决于法律授权、总理与其他副总理的分工以及议会通过的预算。路透社援引分析人士称，这项任命或使总理米舒斯京的部分协调空间缩小。这是对政府内部权力变化的分析，而非可直接测量的事实。外界还需要观察是否出现具体税改、财政支出削减和预算执行结果，才能判断此次改组的真实政策后果。\cite{ReutersRussiaCabinet}
}

\newsitem{中国与塔吉克斯坦加强安全合作：中亚边境议题为何牵动多方外交？}{中国外交与中亚安全｜10 月 19 日路透社与新华社消息}{
十月十九日，中国国务院总理李强在与塔吉克斯坦总统拉赫蒙会谈时提出加强执法、安全和情报交流合作，并希望塔方采取措施保护中国公民。路透社援引新华社报道，中方提出共同打击其所称的恐怖主义、分裂主义和极端主义。这里的“三股势力”是中国政府采用的政策概念，并非对任何具体个人或组织已经作出的司法裁定；涉及具体案件时仍应核查证据和法律程序。\cite{ReutersTajik}

塔吉克斯坦与阿富汗接壤，山区边境地形复杂，走私、跨境武装活动与边境执法一直是当地政府关切的问题。路透社称，去年末发生的袭击造成五名中国公民死亡，相关当局认为袭击者来自阿富汗方向。应注意，“当局认定的袭击来源”与经独立调查确认的完整责任链条之间仍有区别。对在当地参与基础设施、矿业和交通项目的企业而言，人员安全与投资稳定具有直接联系。

这次会晤也发生在大国持续加强中亚外交接触的背景下。路透社报道，美国国务卿鲁比奥同期访问乌兹别克斯坦，中国总理则在访问乌兹别克斯坦后前往塔吉克斯坦。把这种外交活动描述为各国争夺影响力，有一定解释价值，但并不意味着中亚国家没有自主外交目标。塔吉克斯坦还与俄罗斯保持紧密安全联系，俄方在当地设有军事基地。中国、俄罗斯、美国和中亚国家的利益有重叠也有差异；后续可核实的重点是具体安全协议、跨境执法合作范围及投资项目的实际执行情况。\cite{ReutersTajik}
}

\newsitem{资本市场稳定与科技融资：监管层警惕跨市场风险传染}{金融监管｜路透社 10 月 19 日报道}{
中国证监会主席吴清十月十九日在北京金融论坛表示，将加强对国内外市场的监测、风险评估与政策工具储备，特别关注股票、债券、汇率市场之间联系增强所带来的跨境、跨市场和跨行业风险传染。路透社指出，沪深三百指数此前已连续十周下跌，科创五十指数较六月末高点下跌接近百分之四十。两项市场行情反映的是特定时间窗口的价格变化，不代表所有中国上市企业价值同步下降，更不能直接推导实体经济在同一期间萎缩同样幅度。\cite{ReutersCSRC}

监管层还提出推进科创板和创业板改革，完善科技企业识别和筛选机制，并推动首批十八只主动管理型交易所交易基金（ETF）的注册。传统被动指数 ETF 主要追踪预设指数，主动管理型 ETF 则由管理人依投资策略选择证券；二者都面临市场风险，主动管理并不意味着收益有保证。发行新产品可以丰富投资工具，却不能单独解决上市公司盈利、信息披露和投资者信任等问题。

所谓跨市场风险传染，可能从杠杆融资、抵押品价格、汇率变化或金融机构共同持仓等渠道发生。例如，一类资产价格下跌迫使机构出售其他资产补充保证金，即使后者基本面并未变化，也可能承受抛售压力。因此监管层强调关注与境内个股挂钩的离岸杠杆产品，其风险不只在产品本身，还在于跨境交易链条是否透明、保证金规则是否充分以及市场参与者是否了解损失可能超过预期。后续应追踪具体规则文本和实施效果，而非仅凭监管表态判断市场已经企稳。\cite{ReutersCSRC}
}

\newsitem{霍尔木兹海峡安全争议：伊朗与俄罗斯外长会晤释放什么信号？}{中东外交｜10 月 19 日路透社现场报道}{
十月十九日，伊朗外长阿拉格齐与到访的俄罗斯外长拉夫罗夫共同出席记者会。阿拉格齐将霍尔木兹海峡持续紧张归咎于美国，并主张地区问题应通过外交途径解决，提出由周边国家共同构建区域安全安排。路透社记录了上述表态，但报道并未独立确认某一方对海峡全部安全事件负有责任。对外交新闻而言，“某国指责”与“已经查明的责任”必须严格分开。\cite{ReutersIran}

霍尔木兹海峡连接波斯湾与阿曼湾，是全球能源运输的重要航道。航道风险上升可能影响保险费用、油轮航线和能源价格，但价格变化同时受产油国政策、全球需求、库存和金融市场预期影响，不能将所有油价波动都归因于某次记者会。伊朗提出减少域外力量参与的区域安全方案，体现了其外交立场；海湾其他国家可能有不同的安全依赖和威胁判断，因此能否形成共同机制，仍需要看到具体谈判对象、航行自由保障、核查机制与争端解决程序。

这次会晤与俄罗斯在乌克兰战事中的外交处境也有联系，但公开材料不足以证明双方达成了新的海峡军事部署安排。判断局势是否真正缓和，应关注商船通行、袭击与拦截事件的独立记录、各国海军部署和公开签署的协议，而不应把外交辞令当作行动已经发生。\cite{ReutersIran}
}

\newsitem{数学拾趣｜平面图的“四色”不是你以为的那个四色定理}{组合数学预印本｜Sedlar 与 Škrekovski，10 月 19 日新稿列表}{
提到平面图与四色，人们通常想到四色定理：一张平面地图只需四种颜色，使相邻区域颜色不同。十月十九日出现在 arXiv 新稿列表的一篇预印本研究的是另一种着色问题：\emph{闭邻域无冲突着色}。对于图中任意顶点 $v$，把它自己及与它相邻的顶点组成闭邻域 $N[v]$；要求在这个集合里，至少有一种颜色\emph{恰好出现一次}。这里并不要求相邻顶点颜色不同，因此它与普通的图着色不是同一概念。\cite{SedlarColour}

举个最简单的例子：一个中心顶点连接若干叶子，假如中心染蓝、所有叶子染红，那么中心的闭邻域中蓝色只出现一次；每个叶子的闭邻域中蓝色和红色也各出现一次。这说明某些图即使许多相邻顶点同色，仍满足无冲突条件。真正困难的是设计一个对所有顶点都成立的全局着色，并确定最坏情况下究竟需要多少种颜色。

作者 Jelena Sedlar 与 Riste Škrekovski 声称推翻了此前“三色足够”的猜想：他们构造了需要四种颜色的平面图，其中包括一个二十五个顶点的实例，并进一步证明判断平面图是否能用三色满足该条件是 NP 完全问题。NP 完全性涉及一类计算判定问题的复杂性，不等于某个特定二十五顶点图无法被计算机检查。作者还给出无限族构造。论文初稿提交于十月十五日，目前属于预印本；本刊核对了摘要和定义，但未独立检验证明或复杂性归约。它的启发在于：看似只是把“四色定理”换了个说法，量词和邻域条件一变，问题的结构与算法难度就可能完全不同。\cite{SedlarColour}
}

\par\medskip{\large\bfseries\color{newsblue}参考资料与原始出处}\par
\begingroup\footnotesize\renewcommand{\refname}{}\bibliographystyle{unsrt}\bibliography{refs}\endgroup
\end{documents}
\end{document}